% GENERATED FILE. Edit canonical lessons, then run npm run generate:books.
% canonical-source-hash: fnv1a64:dd843adbc5e637d3
% canonical-lessons: SA-C148-napitah, SA-C148-pacakah, SA-C148-malakarah, SA-C148-dhivarah, SA-C148-gopalah

\chapter{Barber, Cook, Gardener, Fisherman, Cowherd}
\label{ch:sa-barber-cook-gardener-fisherman-cowherd}
\hlchaptermodality{148}

\begin{chapteropening}
\textbf{By the end of this chapter:} \emph{I can say a barber, a cook, a gardener, a fisherman and a cowherd.}

Complete the last lesson of chapter 148: I can say a barber, a cook, a gardener, a fisherman and a cowherd.
\end{chapteropening}

\section[\texorpdfstring{\sk{नापितः} (nāpitaḥ)}{नापितः (nāpitaḥ)}]{\sk{नापितः} (nāpitaḥ) — a barber}

\label{lesson:SA-C148-napitah}



\begin{quote}
Before the new one: say the Sanskrit for a minister, then the Sanskrit for a general.
\end{quote}

\subsection*{You'll want to know: \sk{नापितः}}
\textbf{\sk{नापितः}} — \emph{nāpitaḥ} — \textquotedblleft{}a barber\textquotedblright{}.

\begin{grammarlens}[title={What you've built}]
One more word for this chapter.
\end{grammarlens}

\subsection*{Your turn}
\begin{itemize}
\raggedright
  \item \emph{Say it:} \emph{nāpitaḥ}
  \item \emph{Say it:} \emph{nāpitaḥ}, once more
  \item \emph{Say it:} say \emph{nāpitaḥ}
  \item \emph{Recall:} say the Sanskrit for a minister, then the Sanskrit for a general, then say \emph{nāpitaḥ} again
\end{itemize}

\begin{tcolorbox}[breakable,colback=teal!4,colframe=teal!35!black,title={Before you move on}]
What does \sk{नापितः} mean? (\textquotedblleft{}A barber\textquotedblright{}.) Say it once more.
\end{tcolorbox}
\section[\texorpdfstring{\sk{पाचकः} (pācakaḥ)}{पाचकः (pācakaḥ)}]{\sk{पाचकः} (pācakaḥ) — a cook}

\label{lesson:SA-C148-pacakah}



\begin{quote}
Before the new one: say the Sanskrit for a general, then the Sanskrit for a barber.
\end{quote}

\subsection*{You'll want to know: \sk{पाचकः}}
\textbf{\sk{पाचकः}} — \emph{pācakaḥ} — \textquotedblleft{}a cook\textquotedblright{}.

\begin{grammarlens}[title={What you've built}]
One more word for this chapter.
\end{grammarlens}

\subsection*{Your turn}
\begin{itemize}
\raggedright
  \item \emph{Say it:} \emph{pācakaḥ}
  \item \emph{Say it:} \emph{pācakaḥ}, once more
  \item \emph{Say it:} say \emph{pācakaḥ}
  \item \emph{Recall:} say the Sanskrit for a general, then the Sanskrit for a barber, then say \emph{pācakaḥ} again
\end{itemize}

\begin{tcolorbox}[breakable,colback=teal!4,colframe=teal!35!black,title={Before you move on}]
What does \sk{पाचकः} mean? (\textquotedblleft{}A cook\textquotedblright{}.) Say it once more.
\end{tcolorbox}
\section[\texorpdfstring{\sk{मालाकारः} (mālākāraḥ)}{मालाकारः (mālākāraḥ)}]{\sk{मालाकारः} (mālākāraḥ) — a gardener, a garland-maker}

\label{lesson:SA-C148-malakarah}



\begin{quote}
Before the new one: say the Sanskrit for a barber, then the Sanskrit for a cook.
\end{quote}

\subsection*{You'll want to know: \sk{मालाकारः}}
\textbf{\sk{मालाकारः}} — \emph{mālākāraḥ} — \textquotedblleft{}a gardener, a garland-maker\textquotedblright{}.

\begin{grammarlens}[title={What you've built}]
One more word for this chapter.
\end{grammarlens}

\subsection*{Your turn}
\begin{itemize}
\raggedright
  \item \emph{Say it:} \emph{mālākāraḥ}
  \item \emph{Say it:} \emph{mālākāraḥ}, once more
  \item \emph{Say it:} say \emph{mālākāraḥ}
  \item \emph{Recall:} say the Sanskrit for a barber, then the Sanskrit for a cook, then say \emph{mālākāraḥ} again
\end{itemize}

\begin{tcolorbox}[breakable,colback=teal!4,colframe=teal!35!black,title={Before you move on}]
What does \sk{मालाकारः} mean? (\textquotedblleft{}A gardener, a garland-maker\textquotedblright{}.) Say it once more.
\end{tcolorbox}
\section[\texorpdfstring{\sk{धीवरः} (dhīvaraḥ)}{धीवरः (dhīvaraḥ)}]{\sk{धीवरः} (dhīvaraḥ) — a fisherman}

\label{lesson:SA-C148-dhivarah}



\begin{quote}
Before the new one: say the Sanskrit for a cook, then the Sanskrit for a gardener.
\end{quote}

\subsection*{You'll want to know: \sk{धीवरः}}
\textbf{\sk{धीवरः}} — \emph{dhīvaraḥ} — \textquotedblleft{}a fisherman\textquotedblright{}.

\begin{grammarlens}[title={What you've built}]
One more word for this chapter.
\end{grammarlens}

\subsection*{Your turn}
\begin{itemize}
\raggedright
  \item \emph{Say it:} \emph{dhīvaraḥ}
  \item \emph{Say it:} \emph{dhīvaraḥ}, once more
  \item \emph{Say it:} say \emph{dhīvaraḥ}
  \item \emph{Recall:} say the Sanskrit for a cook, then the Sanskrit for a gardener, then say \emph{dhīvaraḥ} again
\end{itemize}

\begin{tcolorbox}[breakable,colback=teal!4,colframe=teal!35!black,title={Before you move on}]
What does \sk{धीवरः} mean? (\textquotedblleft{}A fisherman\textquotedblright{}.) Say it once more.
\end{tcolorbox}
\section[\texorpdfstring{\sk{गोपालः} (gopālaḥ)}{गोपालः (gopālaḥ)}]{\sk{गोपालः} (gopālaḥ) — a cowherd}

\label{lesson:SA-C148-gopalah}



\begin{quote}
Before the new one: say the Sanskrit for a gardener, then the Sanskrit for a fisherman.
\end{quote}

\subsection*{You'll want to know: \sk{गोपालः}}
\textbf{\sk{गोपालः}} — \emph{gopālaḥ} — \textquotedblleft{}a cowherd\textquotedblright{}.

\begin{grammarlens}[title={What you've built}]
One more word for this chapter.
\end{grammarlens}

\subsection*{Your turn}
\begin{itemize}
\raggedright
  \item \emph{Say it:} \emph{gopālaḥ}
  \item \emph{Say it:} \emph{gopālaḥ}, once more
  \item \emph{Say it:} say \emph{gopālaḥ}
  \item \emph{Recall:} say the Sanskrit for a gardener, then the Sanskrit for a fisherman, then say \emph{gopālaḥ} again
\end{itemize}

\begin{tcolorbox}[breakable,colback=teal!4,colframe=teal!35!black,title={Before you move on}]
What does \sk{गोपालः} mean? (\textquotedblleft{}A cowherd\textquotedblright{}.) Say it once more.
\end{tcolorbox}
